% Section 2.3, from lectures/L02/README.md: the objectives, the evaluation questions, and the next
% lecture.

\section{Review}
\label{c2:sec:review}

\needspace{6\baselineskip}
\subsection*{What you should be able to do now}
After this chapter, you should:
\begin{itemize}
  \item Be able to create simple classes in C++.
  \item Understand the purpose of constructors and destructors.
  \item Understand the concept of encapsulation in object-oriented programming.
  \item Understand how classes are organized across header and source files.
  \item Be familiar with copy and move semantics.
  \item Understand the purpose of the keywords \code{explicit}, \code{final}, \code{default}, and
    \code{delete}.
  \item Be able to define and use enumeration classes (\code{enum class}).
\end{itemize}

\needspace{6\baselineskip}
\subsection*{Questions to test yourself}
You should be able to explain:
\begin{enumerate}
  \item What is the difference between a struct and a class in C++?
  \item What does encapsulation mean?
  \item What is a constructor, and when is it called?
  \item What is a destructor, and when is it called?
  \item What is the meaning of the keywords \code{explicit} and \code{final}?
  \item What are \code{default} and \code{delete} used for?
\end{enumerate}

\needspace{6\baselineskip}
\subsection*{Looking ahead}
Chapter~\ref{c3:ch:interfaces} is about inheritance and interfaces. You will learn how larger
systems can be structured using:
\begin{itemize}
  \item Class inheritance.
  \item Abstract interfaces.
  \item Interface-based driver design in embedded systems.
\end{itemize}
